\documentclass[10pt,a4paper]{amsart}
\usepackage{amsmath,amssymb,amsthm,graphicx}
% Retained minimal preamble for the Stacks Project chapter more-morphisms.tex (More on Morphisms).
% Verbatim excerpt of the chapter preamble: Xy-pic setup (native preamble lines 14-18) and the native
% theorem declarations with the mathematical macros (native preamble lines 154-217). Copied mechanically;
% nothing added, removed or redefined. The wrapper class and the amsmath/amssymb/amsthm/graphicx base
% packages are supplied by the assembler, so the class/externaldocument part of the native preamble is omitted.
% For commutative diagrams we use Xy-pic
\usepackage[all]{xy}

% We use 2cell for 2-commutative diagrams.
\xyoption{2cell}
% Theorem environments.
%
\theoremstyle{plain}
\newtheorem{theorem}[subsection]{Theorem}
\newtheorem{proposition}[subsection]{Proposition}
\newtheorem{lemma}[subsection]{Lemma}

\theoremstyle{definition}
\newtheorem{definition}[subsection]{Definition}
\newtheorem{example}[subsection]{Example}
\newtheorem{exercise}[subsection]{Exercise}
\newtheorem{situation}[subsection]{Situation}

\theoremstyle{remark}
\newtheorem{remark}[subsection]{Remark}
\newtheorem{remarks}[subsection]{Remarks}

\numberwithin{equation}{subsection}

% Macros
%
\def\lim{\mathop{\mathrm{lim}}\nolimits}
\def\colim{\mathop{\mathrm{colim}}\nolimits}
\def\Spec{\mathop{\mathrm{Spec}}}
\def\Hom{\mathop{\mathrm{Hom}}\nolimits}
\def\Ext{\mathop{\mathrm{Ext}}\nolimits}
\def\SheafHom{\mathop{\mathcal{H}\!\mathit{om}}\nolimits}
\def\SheafExt{\mathop{\mathcal{E}\!\mathit{xt}}\nolimits}
\def\Sch{\mathit{Sch}}
\def\Mor{\mathop{\mathrm{Mor}}\nolimits}
\def\Ob{\mathop{\mathrm{Ob}}\nolimits}
\def\Sh{\mathop{\mathit{Sh}}\nolimits}
\def\NL{\mathop{N\!L}\nolimits}
\def\CH{\mathop{\mathrm{CH}}\nolimits}
\def\proetale{{pro\text{-}\acute{e}tale}}
\def\etale{{\acute{e}tale}}
\def\QCoh{\mathit{QCoh}}
\def\Ker{\mathop{\mathrm{Ker}}}
\def\Im{\mathop{\mathrm{Im}}}
\def\Coker{\mathop{\mathrm{Coker}}}
\def\Coim{\mathop{\mathrm{Coim}}}

% Boxtimes
%
\DeclareMathSymbol{\boxtimes}{\mathbin}{AMSa}{"02}

%
% Macros for moduli stacks/spaces
%
\def\QCohstack{\mathcal{QC}\!\mathit{oh}}
\def\Cohstack{\mathcal{C}\!\mathit{oh}}
\def\Spacesstack{\mathcal{S}\!\mathit{paces}}
\def\Quotfunctor{\mathrm{Quot}}
\def\Hilbfunctor{\mathrm{Hilb}}
\def\Curvesstack{\mathcal{C}\!\mathit{urves}}
\def\Polarizedstack{\mathcal{P}\!\mathit{olarized}}
\def\Complexesstack{\mathcal{C}\!\mathit{omplexes}}
% \Pic is the operator that assigns to X its picard group, usage \Pic(X)
% \Picardstack_{X/B} denotes the Picard stack of X over B
% \Picardfunctor_{X/B} denotes the Picard functor of X over B
\def\Pic{\mathop{\mathrm{Pic}}\nolimits}
\def\Picardstack{\mathcal{P}\!\mathit{ic}}
\def\Picardfunctor{\mathrm{Pic}}
\def\Deformationcategory{\mathcal{D}\!\mathit{ef}}
\newenvironment{reference}{}{}
\newenvironment{slogan}{}{}
\newenvironment{history}{}{}
\usepackage{cleveref}
\crefname{originallemma}{original lemma}{original lemmas}
\Crefname{originallemma}{Original lemma}{Original lemmas}
\makeatletter
\newlabel{coherent-lemma-flat-base-change-cohomology}{{coherent-lemma-flat-base-change-cohomology}{source}{}{sourcepointer-1}{}}
\newlabel{coherent-lemma-flat-base-change-cohomology@cref}{{[originallemma][0][]coherent-lemma-flat-base-change-cohomology}{[1][0][]source}{}{}{}}
\makeatother
\crefname{originallemma}{original lemma}{original lemmas}
\Crefname{originallemma}{Original lemma}{Original lemmas}
\makeatletter
\newlabel{constructions-lemma-spec-properties}{{constructions-lemma-spec-properties}{source}{}{sourcepointer-2}{}}
\newlabel{constructions-lemma-spec-properties@cref}{{[originallemma][0][]constructions-lemma-spec-properties}{[1][0][]source}{}{}{}}
\makeatother
\crefname{originallemma}{original lemma}{original lemmas}
\Crefname{originallemma}{Original lemma}{Original lemmas}
\makeatletter
\newlabel{lemma-integral-closure-smooth-pullback}{{lemma-integral-closure-smooth-pullback}{source}{}{sourcepointer-3}{}}
\newlabel{lemma-integral-closure-smooth-pullback@cref}{{[originallemma][0][]lemma-integral-closure-smooth-pullback}{[1][0][]source}{}{}{}}
\makeatother
\crefname{originallemma}{original lemma}{original lemmas}
\Crefname{originallemma}{Original lemma}{Original lemmas}
\makeatletter
\newlabel{morphisms-lemma-base-change-finite}{{01WL}{source}{}{sourcepointer-4}{}}
\newlabel{morphisms-lemma-base-change-finite@cref}{{[originallemma][0][]01WL}{[1][0][]source}{}{}{}}
\makeatother
\crefname{originallemma}{original lemma}{original lemmas}
\Crefname{originallemma}{Original lemma}{Original lemmas}
\makeatletter
\newlabel{morphisms-lemma-characterize-normalization}{{morphisms-lemma-characterize-normalization}{source}{}{sourcepointer-5}{}}
\newlabel{morphisms-lemma-characterize-normalization@cref}{{[originallemma][0][]morphisms-lemma-characterize-normalization}{[1][0][]source}{}{}{}}
\makeatother
\title{Normalization of schemes commutes with smooth base change}
\author{The Stacks Project Authors}
\date{stacks-project:more-morphisms.tex:local-snapshot-20261011; GNU FDL 1.2 or later}
\begin{document}
\maketitle
\noindent\textit{Editorial note.} Statement span is lines 5373-5386 of sources/stacks-more-morphisms/pinned-original-source.tex; proof span is lines 5388-5407, the complete author proof (18 body lines).\par\medskip
\noindent\textit{Editorial note.} Navigation stubs: 5 cited-but-not-included labels are supplied as non-load-bearing pointers with their pinned tags where the snapshot tag index has them.\par\medskip
\noindent\textit{Editorial note.} extra\_packages is empty on purpose: the retained preamble already loads xy with options and a second declaration clashes in the assembler wrapper.\par\medskip
\noindent\textit{Editorial note.} Text is retained verbatim, including the author's own typos and placeholders.\par\medskip
\noindent\textit{Allowed standard prerequisites (editorial).}\begin{itemize}
\item Explicit prerequisite boundary. Normalization of a scheme in a quasi-coherent sheaf of algebras, integral closure, and smooth base change; the reduction is retained.
\end{itemize}\medskip
\ifdefined\hypertarget\hypertarget{sourcepointer-1}{}\fi
\noindent\textit{Source reference.} Original lemma coherent-lemma-flat-base-change-cohomology: sources/stacks-more-morphisms/pinned-original-source.tex, source lines 5407--5407. Cited by the retained statement or proof but not part of the collected spans; supplied only as a navigation pointer. The argument does not depend on its content.\par\medskip
\ifdefined\hypertarget\hypertarget{sourcepointer-2}{}\fi
\noindent\textit{Source reference.} Original lemma constructions-lemma-spec-properties: sources/stacks-more-morphisms/pinned-original-source.tex, source lines 5407--5407. Cited by the retained statement or proof but not part of the collected spans; supplied only as a navigation pointer. The argument does not depend on its content.\par\medskip
\ifdefined\hypertarget\hypertarget{sourcepointer-3}{}\fi
\noindent\textit{Source reference.} Original lemma lemma-integral-closure-smooth-pullback: sources/stacks-more-morphisms/pinned-original-source.tex, source lines 5359--5365. Cited by the retained statement or proof but not part of the collected spans; supplied only as a navigation pointer. The argument does not depend on its content.\par\medskip
\ifdefined\hypertarget\hypertarget{sourcepointer-4}{}\fi
\noindent\textit{Source reference.} Original lemma 01WL: sources/stacks-more-morphisms/pinned-original-source.tex, source lines 5407--5407. Cited by the retained statement or proof but not part of the collected spans; supplied only as a navigation pointer. The argument does not depend on its content.\par\medskip
\ifdefined\hypertarget\hypertarget{sourcepointer-5}{}\fi
\noindent\textit{Source reference.} Original lemma morphisms-lemma-characterize-normalization: sources/stacks-more-morphisms/pinned-original-source.tex, source lines 5407--5407. Cited by the retained statement or proof but not part of the collected spans; supplied only as a navigation pointer. The argument does not depend on its content.\par\medskip
\begin{lemma}[Normalization commutes with smooth base change]
\label{lemma-normalization-smooth-localization}
Let
$$
\xymatrix{
Y_2 \ar[r] \ar[d]_{f_2} & Y_1 \ar[d]^{f_1} \\
X_2 \ar[r]^\varphi & X_1
}
$$
be a fibre square in the category of schemes. Assume $f_1$ is quasi-compact
and quasi-separated, and $\varphi$ is smooth.
Let $Y_i \to X_i' \to X_i$ be the normalization of $X_i$ in $Y_i$.
Then $X_2' \cong X_2 \times_{X_1} X_1'$.
\end{lemma}

\begin{proof}
The base change of the factorization $Y_1 \to X_1' \to X_1$ to $X_2$
is a factorization $Y_2 \to X_2 \times_{X_1} X_1' \to X_2$ and
$X_2 \times_{X_1} X_1' \to X_2$ is integral
(Morphisms, Lemma \ref{morphisms-lemma-base-change-finite}).
Hence we get a morphism
$h : X_2' \to X_2 \times_{X_1} X_1'$ by the universal property of
Morphisms, Lemma \ref{morphisms-lemma-characterize-normalization}.
Observe that $X_2'$ is the relative spectrum of the integral closure
of $\mathcal{O}_{X_2}$ in $f_{2, *}\mathcal{O}_{Y_2}$.
If $\mathcal{A}' \subset f_{1, *}\mathcal{O}_{Y_1}$ denotes the integral
closure of $\mathcal{O}_{X_1}$, then $X_2 \times_{X_1} X_1'$ is the
relative spectrum of $\varphi^*\mathcal{A}'$, see
Constructions, Lemma \ref{constructions-lemma-spec-properties}.
By
Cohomology of Schemes, Lemma \ref{coherent-lemma-flat-base-change-cohomology}
we know that $f_{2, *}\mathcal{O}_{Y_2} = \varphi^*f_{1, *}\mathcal{O}_{Y_1}$.
Hence the result follows from
Lemma \ref{lemma-integral-closure-smooth-pullback}.
\end{proof}

\end{document}
